\section{Advanced Hypothesis Testing: Two Means, Proportions, Variances, and Homogeneity}

In the preceding sections of this chapter, we established the philosophical and procedural foundations of statistical hypothesis testing: we formulated null ($H_0$) and alternative ($H_a$) hypotheses, analyzed Type I ($\alpha$) and Type II ($\beta$) errors, and studied classical $Z$ and $t$ tests for a single mean. Furthermore, we introduced the $\chi^2$ test to verify the goodness of fit of a distribution and to test for independence between two categorical variables within a single population.

In the empirical practice of data science, experimental analytics, and industrial design, researchers continually encounter comparative and multidimensional questions: Is algorithm $A$ significantly faster than algorithm $B$? Does the click-through or conversion rate differ across three distinct demographic segments? Does introducing a new sensor increase stability (reduce variance) in a system?

In this section, we fully and rigorously develop the theory and methodology for testing hypotheses concerning two means, one and two proportions, one and two variances, as well as tests of homogeneity and the simultaneous comparison of multiple proportions.

\subsection{Hypothesis Tests for Two Means ($\mu_1 - \mu_2$)}

When comparing the population means of two groups ($\mu_1$ and $\mu_2$), we test a null hypothesis of the form:
\begin{align}
	H_0: \mu_1 - \mu_2 = \Delta_0,
\end{align}
where $\Delta_0$ is the hypothesized difference (in the vast majority of comparative studies, no initial difference is postulated, meaning $\Delta_0 = 0 \implies \mu_1 = \mu_2$).

The alternative hypothesis ($H_a$) can take three forms, determining the structure of the rejection region:
\begin{itemize}
	\item \textbf{Two-tailed test:} $H_a: \mu_1 - \mu_2 \neq \Delta_0$.
	\item \textbf{Left-tailed test:} $H_a: \mu_1 - \mu_2 < \Delta_0$.
	\item \textbf{Right-tailed test:} $H_a: \mu_1 - \mu_2 > \Delta_0$.
\end{itemize}

The test statistic to use depends on whether the population variances are known or unknown, and whether the samples are independent or paired.

\begin{teorema}[Case 1: Independent samples with known population variances]
	Let two independent random samples of sizes $n_1$ and $n_2$ be drawn from normal populations with means $\mu_1, \mu_2$ and known variances $\sigma_1^2, \sigma_2^2$. (If $n_1, n_2 \geq 30$, population normality is not strictly required due to the Central Limit Theorem). Under the null hypothesis $H_0: \mu_1 - \mu_2 = \Delta_0$, the standardized statistic:
	\begin{align}
		Z = \frac{(\bar{X}_1 - \bar{X}_2) - \Delta_0}{\sqrt{\frac{\sigma_1^2}{n_1} + \frac{\sigma_2^2}{n_2}}}
	\end{align}
	follows exactly a standard normal distribution $N(0,1)$. We reject $H_0$ at significance level $\alpha$ if $|Z| > Z_{\alpha/2}$ (two-tailed), $Z < -Z_\alpha$ (left-tailed), or $Z > Z_\alpha$ (right-tailed).
\end{teorema}

\begin{teorema}[Case 2: Independent samples with unknown but equal variances (Homoscedasticity)]
	When $\sigma_1^2, \sigma_2^2$ are unknown but can be reasonably assumed equal ($\sigma_1^2 = \sigma_2^2 = \sigma^2$), we estimate the common variance using the pooled sample variance $S_p^2$:
	\begin{align}
		S_p^2 = \frac{(n_1 - 1)S_1^2 + (n_2 - 1)S_2^2}{n_1 + n_2 - 2}.
	\end{align}
	The test statistic under $H_0: \mu_1 - \mu_2 = \Delta_0$ is:
	\begin{align}
		t = \frac{(\bar{X}_1 - \bar{X}_2) - \Delta_0}{S_p \sqrt{\frac{1}{n_1} + \frac{1}{n_2}}},
	\end{align}
	which follows a Student's $t$-distribution with $\nu = n_1 + n_2 - 2$ degrees of freedom.
\end{teorema}

\begin{teorema}[Case 3: Independent samples with unknown and unequal variances (Welch)]
	If the population variances are unknown and notably different ($\sigma_1^2 \neq \sigma_2^2$), we cannot pool the sample variances. The test statistic is:
	\begin{align}
		t = \frac{(\bar{X}_1 - \bar{X}_2) - \Delta_0}{\sqrt{\frac{S_1^2}{n_1} + \frac{S_2^2}{n_2}}},
	\end{align}
	whose sampling distribution is excellently approximated by a Student's $t$-distribution with Welch-Satterthwaite degrees of freedom:
	\begin{align}
		\nu = \frac{\left( \frac{S_1^2}{n_1} + \frac{S_2^2}{n_2} \right)^2}{\frac{\left(S_1^2/n_1\right)^2}{n_1 - 1} + \frac{\left(S_2^2/n_2\right)^2}{n_2 - 1}}.
	\end{align}
\end{teorema}

\begin{definicion}[Case 4: Paired or dependent samples]
	When measurements from both groups are taken on the same subjects or paired experimental units ($X_{1i}, X_{2i}$), we define the difference variable $D_i = X_{1i} - X_{2i}$. Testing $\mu_1 - \mu_2$ transforms into a one-sample $t$-test on the mean difference $\mu_D$. Under $H_0: \mu_D = \Delta_0$:
	\begin{align}
		t = \frac{\bar{D} - \Delta_0}{S_D / \sqrt{n}},
	\end{align}
	where $\bar{D} = \frac{1}{n}\sum D_i$ and $S_D = \sqrt{\frac{1}{n-1}\sum (D_i - \bar{D})^2}$. The statistic follows a Student's $t$-distribution with $\nu = n - 1$ degrees of freedom.
\end{definicion}

\subsection{Hypothesis Tests Concerning Proportions}

In statistical modeling and machine learning, evaluating categorical or binary variables (success/failure) requires testing population proportions.

\begin{teorema}[$Z$-test for a single population proportion]
	Let $X$ be the number of successes in a random sample of size $n$ from a binomial distribution with true probability $p$. To test the null hypothesis $H_0: p = p_0$ when the sample size satisfies normal approximation conditions ($n p_0 \geq 5$ and $n(1-p_0) \geq 5$), the standardized test statistic is:
	\begin{align}
		Z = \frac{\hat{p} - p_0}{\sqrt{\frac{p_0(1-p_0)}{n}}},
	\end{align}
	where $\hat{p} = X/n$ is the observed sample proportion. Note a crucial theoretical difference from the confidence interval: in the denominator of the $Z$ statistic, we use the expected standard deviation \emph{under the null hypothesis} ($p_0(1-p_0)$), rather than the empirical sample variance ($\hat{p}(1-\hat{p})$).
\end{teorema}

\begin{teorema}[$Z$-test for the difference between two proportions]
	To compare the success rates of two independent populations ($p_1$ and $p_2$), we test the null hypothesis of equality $H_0: p_1 - p_2 = 0$ ($p_1 = p_2 = p$).
	
	Since under the null hypothesis both samples theoretically originate from the same population sharing a common proportion $p$, the most efficient point estimator of this proportion is the weighted average of successes across both samples, known as the \emph{pooled proportion} ($\bar{p}$):
	\begin{align}
		\bar{p} = \frac{X_1 + X_2}{n_1 + n_2} = \frac{n_1 \hat{p}_1 + n_2 \hat{p}_2}{n_1 + n_2}.
		\label{eq:prop_combinada}
	\end{align}
	
	The test statistic under $H_0$ is:
	\begin{align}
		Z = \frac{\hat{p}_1 - \hat{p}_2}{\sqrt{\bar{p}(1-\bar{p})\left(\frac{1}{n_1} + \frac{1}{n_2}\right)}},
	\end{align}
	which is asymptotically distributed as a standard normal $N(0,1)$ when $n_1 \bar{p}, n_1(1-\bar{p}), n_2 \bar{p}, n_2(1-\bar{p}) \geq 5$.
\end{teorema}

\subsection{Hypothesis Tests Concerning Variances}

Testing variances and dispersions is indispensable not only for evaluating uniformity or risk, but also for validating homoscedasticity assumptions in parametric models.

\begin{teorema}[$\chi^2$ test for the variance of a normal population]
	To test whether the variance of a normal population equals a hypothesized value $H_0: \sigma^2 = \sigma_0^2$, we calculate the sample variance $S^2$ from a sample of size $n$. The test statistic is:
	\begin{align}
		\chi^2 = \frac{(n-1)S^2}{\sigma_0^2}.
	\end{align}
	Under $H_0$, this statistic follows a chi-square distribution with $\nu = n - 1$ degrees of freedom.
	For a two-tailed test ($H_a: \sigma^2 \neq \sigma_0^2$), we reject $H_0$ at level $\alpha$ if $\chi^2 < \chi^2_{1-\alpha/2, n-1}$ or if $\chi^2 > \chi^2_{\alpha/2, n-1}$.
\end{teorema}

\begin{teorema}[$F$-test for comparing two population variances]
	To test the hypothesis of equal variances between two independent normal populations ($H_0: \sigma_1^2 = \sigma_2^2$), we draw samples of sizes $n_1$ and $n_2$ and calculate their sample variances $S_1^2$ and $S_2^2$. The test statistic is the ratio of the sample variances:
	\begin{align}
		F = \frac{S_1^2}{S_2^2}.
	\end{align}
	Under $H_0$, this statistic follows an $F$-distribution of Snedecor with $\nu_1 = n_1 - 1$ numerator degrees of freedom and $\nu_2 = n_2 - 1$ denominator degrees of freedom.
	For a two-tailed test ($H_a: \sigma_1^2 \neq \sigma_2^2$), $H_0$ is rejected if $F > F_{\alpha/2, \nu_1, \nu_2}$ or if $F < F_{1-\alpha/2, \nu_1, \nu_2} = \frac{1}{F_{\alpha/2, \nu_2, \nu_1}}$.
\end{teorema}

\subsection{Tests of Homogeneity vs. Tests of Independence}

In section \ref{sec:3.9}, we explored the application of the chi-square distribution to contingency tables for testing \emph{independence} between two factors. There is another test of major importance for multidimensional comparisons: the \emph{test of homogeneity}. Although the computational formula ($\sum \frac{(O-E)^2}{E}$) is identical in both cases, there is a deep conceptual difference in experimental design and sampling.

\begin{definicion}[Conceptual Difference Between Independence and Homogeneity]
	\begin{itemize}
		\item \textbf{Test of Independence:} \textbf{A single random sample} of total size $N$ is drawn from one population and subsequently cross-classified by two categorical variables $A$ (with $r$ levels) and $B$ (with $c$ levels). No marginal row or column total is fixed in advance by the researcher (except the grand total $N$). The tested hypothesis is $H_0: P(A_i \cap B_j) = P(A_i)P(B_j)$.
		\item \textbf{Test of Homogeneity:} \textbf{$c$ independent random samples} of predetermined sizes ($n_1, n_2, \ldots, n_c$, fixing column totals by experimental design) are drawn from $c$ distinct subpopulations or groups. For each subpopulation, the distribution of \textbf{a single categorical variable} with $r$ levels is measured. The tested hypothesis is that all $c$ populations share identical categorical distributions:
		\begin{align}
			H_0: p_{i1} = p_{i2} = \ldots = p_{ic} \quad \text{for each category } i = 1, 2, \ldots, r,
		\end{align}
		where $p_{ij}$ is the probability that an observation from subpopulation $j$ falls into category $i$.
	\end{itemize}
\end{definicion}

\begin{teorema}[Chi-square statistic for homogeneity]
	Given a data table with $r$ rows (categories) and $c$ columns (independent populations or sub-groups), if $O_{ij}$ is the observed frequency in cell $(i,j)$, the expected frequency under the null hypothesis of homogeneity is:
	\begin{align}
		E_{ij} = \frac{R_i \cdot C_j}{N},
	\end{align}
	where $R_i = \sum_{j=1}^c O_{ij}$ is the row $i$ total, $C_j = n_j = \sum_{i=1}^r O_{ij}$ is the sample size of population $j$, and $N = \sum C_j$ is the grand total.
	
	The test statistic:
	\begin{align}
		\chi^2 = \sum_{i=1}^r \sum_{j=1}^c \frac{(O_{ij} - E_{ij})^2}{E_{ij}}
	\end{align}
	is asymptotically distributed (when all $E_{ij} \geq 5$) as a chi-square with $\nu = (r-1)(c-1)$ degrees of freedom.
\end{teorema}

\subsection{Tests of Multiple Proportions and Marascuilo's Procedure}

Comparing multiple binomial proportions is a technical special case of the homogeneity test where the dependent variable is dichotomous ($r=2$: Success vs. Failure) across $c = k$ independent populations.

\begin{definicion}[Testing $k$ proportions]
	Let $p_1, p_2, \ldots, p_k$ be the true success proportions in $k$ independent populations. The null hypothesis is:
	\begin{align}
		H_0: p_1 = p_2 = \ldots = p_k = p,
	\end{align}
	versus the alternative hypothesis $H_a$: at least one proportion differs from the rest.
	
	The homogeneity $\chi^2$ statistic applied to the $2 \times k$ table has $\nu = (2-1)(k-1) = k-1$ degrees of freedom. It can be algebraically demonstrated that the statistic can be compactly expressed as:
	\begin{align}
		\chi^2 = \sum_{j=1}^k \frac{(X_j - n_j \bar{p})^2}{n_j \bar{p}(1-\bar{p})},
		\label{eq:chi2_k_prop}
	\end{align}
	where $X_j$ is the number of successes in sample $j$ (of size $n_j$), and $\bar{p} = \frac{\sum X_j}{\sum n_j}$ is the overall pooled proportion of successes across the $k$ samples.
\end{definicion}

When the $\chi^2$ statistic is significant and we reject $H_0$, we know with statistical certainty that the $k$ proportions are not all equal, but the omnibus test does not tell us \emph{which} specific pairs differ. To identify significant pairs without inflating the overall family-wise Type I error rate across $\binom{k}{2}$ individual pairwise proportion difference tests, we use Marascuilo's multiple comparison procedure.

\begin{teorema}[Marascuilo post-hoc comparison procedure]
	If the omnibus $\chi^2$ test of $k$ proportions rejects $H_0$ at significance level $\alpha$, the observed absolute differences between every pair of sample proportions ($|\hat{p}_i - \hat{p}_j|$ for all $1 \leq i < j \leq k$) are compared against a pairwise critical range $r_{ij}$ given by:
	\begin{align}
		r_{ij} = \sqrt{\chi^2_{\alpha, \, k-1}} \sqrt{\frac{\hat{p}_i(1-\hat{p}_i)}{n_i} + \frac{\hat{p}_j(1-\hat{p}_j)}{n_j}},
		\label{eq:marascuilo}
	\end{align}
	where $\chi^2_{\alpha, k-1}$ is the upper critical quantile of the omnibus statistic with $k-1$ degrees of freedom. 
	
	We declare that population proportions $p_i$ and $p_j$ differ significantly at family-wise error rate $\alpha$ if and only if $|\hat{p}_i - \hat{p}_j| > r_{ij}$.
\end{teorema}
